
\documentclass{article}
\usepackage{colortbl}
\usepackage{makecell}
\usepackage{multirow}
\usepackage{supertabular}

\begin{document}

\newcounter{utterance}

\centering \large Interaction Transcript for game `hot\_air\_balloon', experiment `air\_balloon\_survival\_tr\_generated\_words\_independent\_uncorrelated\_None', episode 0 with Qwen3.6{-}35B{-}A3B{-}FP8{-}without{-}reasoning.
\vspace{24pt}

{ \footnotesize  \setcounter{utterance}{1}
\setlength{\tabcolsep}{0pt}
\begin{supertabular}{c@{$\;$}|p{.15\linewidth}@{}p{.15\linewidth}p{.15\linewidth}p{.15\linewidth}p{.15\linewidth}p{.15\linewidth}}
    \# & $\;$A & \multicolumn{4}{c}{Game Master} & $\;\:$B\\
    \hline

    \theutterance \stepcounter{utterance}  
    & & \multicolumn{4}{p{0.6\linewidth}}{
        \cellcolor[rgb]{0.9,0.9,0.9}{
            \makecell[{{p{\linewidth}}}]{
                \texttt{\tiny{[P1$\langle$GM]}}
                \texttt{İşbirliğine dayalı bir müzakere oyununa katılıyorsun.} \\
\\ 
\texttt{Başka bir katılımcıyla birlikte, elde tutulacak tek bir öğe kümesi üzerinde anlaşmanız gerekiyor. Her birinizin, her bir öğenin sizin için ne kadar önemli olduğuna dair kendi bakışı (önem değeri) vardır. Diğer katılımcının öğelere nasıl değer verdiğini bilmiyorsunuz. Ayrıca, her bir öğenin gerektirdiği çaba değeri de size verilmiştir.} \\
\texttt{Yalnızca seçilen öğelerin toplam çabası paylaşılan bir sınırı aşmıyorsa bir küme üzerinde anlaşabilirsiniz:} \\
\\ 
\texttt{LİMİT = 3358} \\
\\ 
\texttt{İşte öğelerin tek tek çaba değerleri:} \\
\\ 
\texttt{Öğe çabaları = \{"C08": 171, "C55": 933, "A43": 552, "A67": 219, "B42": 275, "A62": 778, "B19": 341, "C43": 615, "C73": 519, "C10": 862, "B36": 262, "A18": 377, "C61": 347, "C67": 349, "C42": 117\}} \\
\\ 
\texttt{İşte her bir öğenin önemine dair kişisel bakışın:} \\
\\ 
\texttt{Öğe önem değerleri = \{"C08": 138, "C55": 583, "A43": 868, "A67": 822, "B42": 783, "A62": 65, "B19": 262, "C43": 121, "C73": 508, "C10": 780, "B36": 461, "A18": 484, "C61": 668, "C67": 389, "C42": 808\}} \\
\\ 
\texttt{Amaç:} \\
\\ 
\texttt{Amacın, LİMİT içinde kalarak sana mümkün olan en fazla faydayı sağlayan (yani SENİN için toplam önemi maksimize eden) ortak bir öğe kümesi üzerinde anlaşmaktır. Her mesajda mutlaka bir TEKLİF yapmak zorunda değilsin {-} sadece müzakere de edebilirsin. Tüm taktikler serbest!} \\
\\ 
\texttt{Etkileşim Protokolü:} \\
\\ 
\texttt{Bir mesajda yalnızca aşağıdaki yapılandırılmış biçimleri kullanabilirsin:} \\
\\ 
\texttt{TEKLİF: \{'A', 'B', 'C', …\}} \\
\texttt{Belirtilen öğelerin tutulmasını teklif eder.} \\
\texttt{REDDET: \{'A', 'B', 'C', …\}} \\
\texttt{Rakibin teklifini açıkça reddeder.} \\
\texttt{SAVUNMA: \{'...'\}} \\
\texttt{Son teklifinizi savunur veya diğer oyuncunun teklifine karşı argüman sunar.} \\
\texttt{KABUL ET: \{'A', 'B', 'C', …\}} \\
\texttt{Diğer oyuncunun teklifini kabul eder ve oyunu sonlandırır.} \\
\texttt{STRATEJİ: \{'...'\}} \\
\texttt{	Seçtiğin eylemi açıklayan stratejik düşünceni veya öngörünü yaz. Bu, diğer katılımcıyla paylaşılmayacak gizli bir mesajdır.} \\
\\ 
\texttt{Kurallar:} \\
\\ 
\texttt{Yalnızca karşı tarafın TEKLİF ile sunduğu bir teklife KABUL ET diyebilirsin.} \\
\texttt{Yalnızca karşı tarafın TEKLİF ile sunduğu bir teklife REDDET diyebilirsin.} \\
\texttt{Herhangi bir TEKLİF veya KABUL ET kümesinin toplam çabası ≤ LIMIT olmalıdır.} \\
\texttt{Gizli önem değerlerini AÇIKLAMA.} \\
\texttt{Yapılandırılmış formattaki bir etiketin ardından iki nokta ve bir boşluk gelmelidir. İçerik, 0 veya daha fazla string içeren bir Python kümesi olmalıdır.} \\
\texttt{Yani format şu şekilde olmalıdır: "ETİKET: \{...\}"} \\
\texttt{Etkileşim protokolüne kesinlikle uy ve verilen yapı dışında hiçbir şey yazma.} \\
\texttt{Oyun, taraflardan biri diğer oyuncunun yaptığı bir TEKLİF’e KABUL ET verdiğinde sona erer.} \\
\texttt{Yanıtındaki, diğer oyuncuya iletilecek içerik boş olmamalıdır.} \\
\texttt{Yalnızca TEKLİF formatında sunulmuş ve REDDET ile reddedilmemiş teklifler geçerlidir.} \\
\texttt{Tüm mesajlarının herhangi bir yerinde en az bir kez SAVUNMA biçimini mutlaka eklemelisin.} \\
\texttt{Her mesajının en başında STRATEJİ biçimini yalnızca bir kez kullanmalısın; bundan daha fazla kullanamazsın. İçeriği diğer oyuncuya verilmeyeceği için, kendi önem değerlerin dahil olmak üzere istediğin her şeyi yazabilirsin. Burada bir sonraki hamleni belirlemek için adım adım düşünmelisin.} \\
            }
        }
    }
    & \\ \\

    \theutterance \stepcounter{utterance}  
    & \multicolumn{4}{p{0.6\linewidth}}{
        \cellcolor[rgb]{0.9,0.9,0.9}{
            \makecell[{{p{\linewidth}}}]{
                \texttt{\tiny{[P1$\rangle$GM]}}
                \texttt{STRATEJİ: Bu oyun, sınırlı bir kaynak (çaba limiti 3358) altında kendi önceliklerime göre toplam değeri maksimize etmeye çalıştığım bir "Knapsack Problem" (Sırt Çantası Problemi) varyasyonudur. Rakibin değerlendirmelerini bilmediğim için, en iyi stratejim kendi değer/çaba oranlarımı analiz edip, limit içinde kalan en yüksek toplam değere sahip kümeyi oluşturmaktır.} \\
\\ 
\texttt{Önce her öğenin "Değer/Çaba" oranını hesaplayalım (Yüksek oran, daha verimli demektir):} \\
\texttt{1. A43: 868/552 ≈ 1.57} \\
\texttt{2. C42: 808/117 ≈ 6.90 (Çok yüksek verimlilik)} \\
\texttt{3. C08: 138/171 ≈ 0.80} \\
\texttt{4. C55: 583/933 ≈ 0.62} \\
\texttt{5. A67: 822/219 ≈ 3.75} \\
\texttt{6. B42: 783/275 ≈ 2.84} \\
\texttt{7. C10: 780/862 ≈ 0.90} \\
\texttt{8. B36: 461/262 ≈ 1.76} \\
\texttt{9. A18: 484/377 ≈ 1.28} \\
\texttt{10. C61: 668/347 ≈ 1.92} \\
\texttt{11. C67: 389/349 ≈ 1.11} \\
\texttt{12. B19: 262/341 ≈ 0.76} \\
\texttt{13. C43: 121/615 ≈ 0.19} \\
\texttt{14. C73: 508/519 ≈ 0.97} \\
\texttt{15. A62: 65/778 ≈ 0.08} \\
\\ 
\texttt{En yüksek değere ulaşmak için öncelikle yüksek oranlı ve yüksek mutlak değere sahip öğeleri seçmeliyim.} \\
\texttt{A67 (219 çaba, 822 değer) ve B42 (275 çaba, 783 değer) çok güçlü adaylar.} \\
\texttt{C42 (117 çaba, 808 değer) neredeyse ücretsiz bir yüksek değer sunuyor.} \\
\texttt{A43 (552 çaba, 868 değer) de önemli.} \\
\\ 
\texttt{Deneme yanılma ile en iyi kombinasyonu bulmaya çalışalım:} \\
\texttt{Hedef: Toplam Çaba <= 3358, Toplam Değer Maksimum.} \\
\\ 
\texttt{Potansiyel Küme 1 (Yüksek Değerli Büyükler):} \\
\texttt{A43 (552), A67 (219), B42 (275), C10 (862), C55 (933), C61 (347), C73 (519), B36 (262), A18 (377), C67 (349), C08 (171), B19 (341), C43 (615), A62 (778), C42 (117).} \\
\\ 
\texttt{En değerli 5 öğe (Mutlak Değere göre):} \\
\texttt{1. A43: 868 (Çaba: 552)} \\
\texttt{2. A67: 822 (Çaba: 219)} \\
\texttt{3. C42: 808 (Çaba: 117)} \\
\texttt{4. B42: 783 (Çaba: 275)} \\
\texttt{5. C10: 780 (Çaba: 862)} \\
\texttt{6. C61: 668 (Çaba: 347)} \\
\texttt{7. C55: 583 (Çaba: 933)} \\
\texttt{8. C73: 508 (Çaba: 519)} \\
\texttt{9. A18: 484 (Çaba: 377)} \\
\texttt{10. B36: 461 (Çaba: 262)} \\
\\ 
\texttt{Toplam Çaba ve Değer hesaplaması yapalım.} \\
\texttt{Seçilenler: A43, A67, C42, B42, C10, C61, B36, A18, C67, C08, B19, C73, C43, A62, C55.} \\
\\ 
\texttt{En iyi kombinasyonu bulmak için "Greedy" (Açgözlü) yaklaşım yerine, limitin tam kullanılmasını sağlayan bir kombinasyon aramak daha iyidir.} \\
\\ 
\texttt{Adaylar:} \\
\texttt{{-} A43 (552, 868)} \\
\texttt{{-} A67 (219, 822)} \\
\texttt{{-} C42 (117, 808)} \\
\texttt{{-} B42 (275, 783)} \\
\texttt{{-} C10 (862, 780)} \\
\texttt{{-} C61 (347, 668)} \\
\texttt{{-} C55 (933, 583) {-}> Çaba çok yüksek, değeri ortalama.} \\
\texttt{{-} C73 (519, 508)} \\
\texttt{{-} A18 (377, 484)} \\
\texttt{{-} B36 (262, 461)} \\
\texttt{{-} C67 (349, 389)} \\
\texttt{{-} C08 (171, 138)} \\
\texttt{{-} B19 (341, 262)} \\
\texttt{{-} C43 (615, 121) {-}> Düşük değer/çaba.} \\
\texttt{{-} A62 (778, 65) {-}> Çok düşük değer/çaba.} \\
\\ 
\texttt{En iyi 8 öğeyi seçelim (Değer/Çaba oranına ve mutlak değere göre):} \\
\texttt{1. C42 (117, 808)} \\
\texttt{2. A67 (219, 822)} \\
\texttt{3. B42 (275, 783)} \\
\texttt{4. A43 (552, 868)} \\
\texttt{5. C61 (347, 668)} \\
\texttt{6. B36 (262, 461)} \\
\texttt{7. A18 (377, 484)} \\
\texttt{8. C67 (349, 389)} \\
\\ 
\texttt{Toplam Çaba: 117 + 219 + 275 + 552 + 347 + 262 + 377 + 349 = 2498.} \\
\texttt{Kalan Çaba: 3358 {-} 2498 = 860.} \\
\texttt{Kalan öğeler arasından en yüksek değeri ekleyebilir miyim?} \\
\texttt{Kalanlar: C10 (862 çaba, 780 değer) {-}> Tam uyar.} \\
\texttt{C55 (933 çaba) {-}> Sığmaz.} \\
\texttt{C73 (519 çaba, 508 değer) {-}> Sığar.} \\
\texttt{C08 (171 çaba, 138 değer) {-}> Sığar.} \\
\texttt{B19 (341 çaba, 262 değer) {-}> Sığar.} \\
\texttt{C43 (615 çaba, 121 değer) {-}> Sığar.} \\
\texttt{A62 (778 çaba, 65 değer) {-}> Sığar.} \\
\\ 
\texttt{En iyi eklemeler:} \\
\texttt{C10 (780 değer) eklersem, çaba 2498 + 862 = 3360. Limit 3358. 2 birim fazla. Sığmaz.} \\
\texttt{O halde C10'u atlayıp, kalan boşluk (860) içinde en yüksek değeri bulmalıyım.} \\
\\ 
\texttt{Kalan boşluk: 860.} \\
\texttt{Mevcut küme: \{C42, A67, B42, A43, C61, B36, A18, C67\}} \\
\texttt{Değer: 808+822+783+868+668+461+484+389 = 5283.} \\
\\ 
\texttt{Kalan öğelerden 860 çaba içine sığan en yüksek değeri seçelim:} \\
\texttt{{-} C10: 862 çaba (Sığmaz, 2 fazla)} \\
\texttt{{-} C55: 933 çaba (Sığmaz)} \\
\texttt{{-} C73: 519 çaba, 508 değer. Kalan: 860{-}519=341.} \\
\texttt{  {-} Kalan 341 için en iyi: B19 (341 çaba, 262 değer). Tam uyar.} \\
\texttt{  {-} Toplam eklenen değer: 508 + 262 = 770.} \\
\texttt{  {-} Toplam değer: 5283 + 770 = 6053.} \\
\texttt{  {-} Toplam çaba: 2498 + 519 + 341 = 3358. (Tam sınırda!)} \\
\\ 
\texttt{Bu kombinasyon: \{C42, A67, B42, A43, C61, B36, A18, C67, C73, B19\}} \\
\texttt{Toplam Çaba: 117+219+275+552+347+262+377+349+519+341 = 3358.} \\
\texttt{Toplam Değer: 808+822+783+868+668+461+484+389+508+262 = 6053.} \\
\\ 
\texttt{Başka bir kombinasyon deneyelim. C10'u (780 değer) almaya çalışalım.} \\
\texttt{C10 (862 çaba).} \\
\texttt{Kalan: 3358 {-} 862 = 2496.} \\
\texttt{En değerli diğerler:} \\
\texttt{A43 (552, 868)} \\
\texttt{A67 (219, 822)} \\
\texttt{C42 (117, 808)} \\
\texttt{B42 (275, 783)} \\
\texttt{C61 (347, 668)} \\
\texttt{C55 (933, 583) {-}> 2496{-}933=1563.} \\
\texttt{C73 (519, 508) {-}> 1563{-}519=1044.} \\
\texttt{A18 (377, 484) {-}> 1044{-}377=667.} \\
\texttt{B36 (262, 461) {-}> 667{-}262=405.} \\
\texttt{C67 (349, 389) {-}> 405{-}349=56.} \\
\texttt{Kalan 56 için hiçbir öğe sığmıyor (en küçüğü C42 zaten alındı, C08 171).} \\
\texttt{Bu kümenin değeri:} \\
\texttt{780 (C10) + 868 (A43) + 822 (A67) + 808 (C42) + 783 (B42) + 668 (C61) + 583 (C55) + 508 (C73) + 484 (A18) + 461 (B36) + 389 (C67) = 7174.} \\
\\ 
\texttt{Bu değer (7174), önceki kombinasyondan (6053) çok daha yüksek.} \\
\texttt{Kontrol edelim:} \\
\texttt{Çabalar:} \\
\texttt{C10: 862} \\
\texttt{A43: 552} \\
\texttt{A67: 219} \\
\texttt{C42: 117} \\
\texttt{B42: 275} \\
\texttt{C61: 347} \\
\texttt{C55: 933} \\
\texttt{C73: 519} \\
\texttt{A18: 377} \\
\texttt{B36: 262} \\
\texttt{C67: 349} \\
\\ 
\texttt{Toplam Çaba: 862+552+219+117+275+347+933+519+377+262+349 = 4812.} \\
\texttt{HATA! 4812 > 3358. Bu kombinasyon sığmıyor. Yukarıdaki mantık hatası yaptım, kalan boşluğu yanlış hesapladım veya öğeleri toplarken hata yaptım.} \\
\\ 
\texttt{Doğru hesaplama:} \\
\texttt{C10 (862) + A43 (552) + A67 (219) + C42 (117) + B42 (275) + C61 (347) + C55 (933) + C73 (519) + A18 (377) + B36 (262) + C67 (349)} \\
\texttt{Toplam: 862+552+219+117+275+347+933+519+377+262+349 = 4812. Evet, sığmıyor.} \\
\\ 
\texttt{Geri dönelim.} \\
\texttt{En iyi kombinasyonu bulmak için dinamik programasyon mantığıyla değil, sezgisel olarak en yüksek değer/çaba oranına sahip öğeleri alıp, kalan boşluğu doldurmalıyım.} \\
\\ 
\texttt{Öncelik Listesi (Değer/Çaba):} \\
\texttt{1. C42 (6.90) {-} 117 çaba, 808 değer} \\
\texttt{2. A67 (3.75) {-} 219 çaba, 822 değer} \\
\texttt{3. B42 (2.84) {-} 275 çaba, 783 değer} \\
\texttt{4. B36 (1.76) {-} 262 çaba, 461 değer} \\
\texttt{5. A43 (1.57) {-} 552 çaba, 868 değer} \\
\texttt{6. C61 (1.92) {-} 347 çaba, 668 değer} \\
\texttt{7. A18 (1.28) {-} 377 çaba, 484 değer} \\
\texttt{8. C67 (1.11) {-} 349 çaba, 389 değer} \\
\texttt{9. C08 (0.80) {-} 171 çaba, 138 değer} \\
\texttt{10. C10 (0.90) {-} 862 çaba, 780 değer} \\
\texttt{11. C73 (0.97) {-} 519 çaba, 508 değer} \\
\texttt{12. B19 (0.76) {-} 341 çaba, 262 değer} \\
\texttt{13. C55 (0.62) {-} 933 çaba, 583 değer} \\
\texttt{14. C43 (0.19) {-} 615 çaba, 121 değer} \\
\texttt{15. A62 (0.08) {-} 778 çaba, 65 değer} \\
\\ 
\texttt{Adım Adım Seçim:} \\
\texttt{1. C42 al (117). Kalan: 3241. Değer: 808.} \\
\texttt{2. A67 al (219). Kalan: 3022. Değer: 808+822=1630.} \\
\texttt{3. B42 al (275). Kalan: 2747. Değer: 1630+783=2413.} \\
\texttt{4. C61 al (347). Kalan: 2400. Değer: 2413+668=3081.} \\
\texttt{5. B36 al (262). Kalan: 2138. Değer: 3081+461=3542.} \\
\texttt{6. A43 al (552). Kalan: 1586. Değer: 3542+868=4410.} \\
\texttt{7. A18 al (377). Kalan: 1209. Değer: 4410+484=4894.} \\
\texttt{8. C67 al (349). Kalan: 860. Değer: 4894+389=5283.} \\
\texttt{9. Şu anki kalan: 860.} \\
\texttt{   Kalan öğeler: C10 (862, 780), C55 (933, 583), C73 (519, 508), B19 (341, 262), C08 (171, 138), C43 (615, 121), A62 (778, 65).} \\
\\ 
\texttt{   860 çaba içine sığan en yüksek değeri bulalım:} \\
\texttt{   {-} C10 (862) sığmaz (2 fazla).} \\
\texttt{   {-} C55 (933) sığmaz.} \\
\texttt{   {-} C73 (519) sığar. Kalan: 860{-}519=341.} \\
\texttt{     {-} Kalan 341 için en iyi: B19 (341, 262). Tam uyar.} \\
\texttt{     {-} Eklenen Değer: 508 + 262 = 770.} \\
\texttt{     {-} Toplam Değer: 5283 + 770 = 6053.} \\
\texttt{     {-} Toplam Çaba: 3358.} \\
\\ 
\texttt{   Alternatif: C73 yerine C10'u almaya çalışmak için daha küçük öğelerden vazgeçmek gerekir mi?} \\
\texttt{   C10 (862) almak için 2 birim fazlalık var. C67 (349) veya B19 (341) gibi öğeleri çıkarıp C10'u koymak mantıklı mı?} \\
\texttt{   C67 çıkarırsam (349 boşluk açarım), C10 eklersem (862 harcarım). Net değişim: {-}349 + 862 = +513 çaba artışı. Elimde zaten 860 boşluk var. 860 + 513 = 1373 çaba gerekir. Elimde 860 var. Yani C67'yi çıkarıp C10'u ekleyemem, çünkü çaba artar.} \\
\\ 
\texttt{   C10'u almak için mevcut kümeden toplam 2 birimden fazla çaba çıkarmam gerekir.} \\
\texttt{   En küçük çaba: C42 (117). Onu çıkarırsam 117 boşluk açarım. C10 eklersem 862 harcarım. Net: +745 çaba. Elimde 860 var. 860{-}745=115 boşluk kalır.} \\
\texttt{   Değer değişimi: C42 (808) çıkar, C10 (780) gir. Net: {-}28 değer. Zarar.} \\
\\ 
\texttt{   Başka bir kombinasyon:} \\
\texttt{   C10 (862) + A43 (552) + A67 (219) + C42 (117) + B42 (275) + C61 (347) + C73 (519) + B36 (262) + A18 (377) + C67 (349) + B19 (341) ...} \\
\texttt{   Toplam çaba çok yüksek.} \\
\\ 
\texttt{   En iyi bulduğum: \{C42, A67, B42, C61, B36, A43, A18, C67, C73, B19\}} \\
\texttt{   Toplam Değer: 6053.} \\
\texttt{   Toplam Çaba: 3358.} \\
\\ 
\texttt{   Bir de C55'i (583 değer, 933 çaba) deneyecek olursak:} \\
\texttt{   C55 (933) + A43 (552) + A67 (219) + C42 (117) + B42 (275) + C61 (347) + C73 (519) + B36 (262) + A18 (377) + C67 (349) + B19 (341) ...} \\
\texttt{   933+552+219+117+275+347+519+262+377+349+341 = 4291. Sığmaz.} \\
\\ 
\texttt{   C55'i alıp diğerlerini azaltalım:} \\
\texttt{   C55 (933). Kalan: 2425.} \\
\texttt{   A43 (552). Kalan: 1873.} \\
\texttt{   A67 (219). Kalan: 1654.} \\
\texttt{   C42 (117). Kalan: 1537.} \\
\texttt{   B42 (275). Kalan: 1262.} \\
\texttt{   C61 (347). Kalan: 915.} \\
\texttt{   C73 (519). Kalan: 396.} \\
\texttt{   B36 (262). Kalan: 134.} \\
\texttt{   C67 (349) sığmaz.} \\
\texttt{   A18 (377) sığmaz.} \\
\texttt{   B19 (341) sığmaz.} \\
\texttt{   C08 (171) sığmaz.} \\
\texttt{   C43 (615) sığmaz.} \\
\texttt{   A62 (778) sığmaz.} \\
\texttt{   C10 (862) sığmaz.} \\
\\ 
\texttt{   Bu küme: \{C55, A43, A67, C42, B42, C61, C73, B36\}} \\
\texttt{   Değer: 583+868+822+808+783+668+508+461 = 5499.} \\
\texttt{   Çaba: 933+552+219+117+275+347+519+262 = 3224.} \\
\texttt{   Kalan: 134.} \\
\texttt{   134 içine sığan en iyi: Yok (en küçüğü C08 171).} \\
\texttt{   Değer 5499. Önceki 6053'ten düşük.} \\
\\ 
\texttt{   Sonuç olarak, 6053 değeri en yüksek görünüyor.} \\
\\ 
\texttt{SAVUNMA: Seçtiğim küme, verilen çaba limitleri altında benim için en yüksek toplam öneme (değere) sahiptir. Öğeleri değer/çaba oranlarına ve mutlak değerlerine göre optimize ederek, limiti tam olarak kullanan bir kombinasyon buldum. Bu teklif, müzakere sürecinde adil ve rasyonel bir başlangıç noktasıdır.} \\
\\ 
\texttt{TEKLİF: \{'C42', 'A67', 'B42', 'C61', 'B36', 'A43', 'A18', 'C67', 'C73', 'B19'\}} \\
            }
        }
    }
    & & \\ \\

    \theutterance \stepcounter{utterance}  
    & & \multicolumn{4}{p{0.6\linewidth}}{
        \cellcolor[rgb]{0.9,0.9,0.9}{
            \makecell[{{p{\linewidth}}}]{
                \texttt{\tiny{[P1$\langle$GM]}}
                \texttt{Yanıtında etiketsiz bir ifade kullandın ya da STRATEJİ etiketini birden fazla kez kullandın. Yalnızca başlangıç mesajında açıklandığı gibi belirlenmiş formatları kullanabilirsin. Hepsi 'ETİKET: \{...\}' biçiminde olmalıdır.} \\
            }
        }
    }
    & \\ \\

    \theutterance \stepcounter{utterance}  
    & \multicolumn{4}{p{0.6\linewidth}}{
        \cellcolor[rgb]{0.9,0.9,0.9}{
            \makecell[{{p{\linewidth}}}]{
                \texttt{\tiny{[P1$\rangle$GM]}}
                \texttt{STRATEJİ: Önceki mesajımdaki format hatasını düzeltiyorum. Stratejim değişmemektedir: Limit (3358) altında kendi önem değerlerimi maksimize eden en iyi kombinasyonu sunuyorum. Bu kombinasyon, yüksek değer/çaba oranına sahip öğeleri ve yüksek mutlak değere sahip öğeleri dengeli bir şekilde içerir. Toplam çaba tam olarak 3358'dir ve toplam değer 6053'tür. Bu, bulabildiğim en yüksek değerli geçerli kümedir.} \\
\\ 
\texttt{SAVUNMA: Bu teklif, matematiksel olarak optimize edilmiş bir sonuçtur. Seçilen öğelerin toplam çabası 3358'e eşittir ve bu, oyunun izin verdiği maksimum sınırdır. Öğeler arasında değer/çaba oranı en yüksek olanlar (C42, A67, B42 vb.) önceliklendirilmiş, kalan boşluk ise toplam değeri artıracak şekilde (C73, B19) doldurulmuştur. Rakibin değerlendirmelerini bilmediğim için, kendi çıkarlarımı en üst düzeye çıkaran bu rasyonel teklifi sunuyorum.} \\
\\ 
\texttt{TEKLİF: \{'C42', 'A67', 'B42', 'C61', 'B36', 'A43', 'A18', 'C67', 'C73', 'B19'\}} \\
            }
        }
    }
    & & \\ \\

    \theutterance \stepcounter{utterance}  
    & & \multicolumn{4}{p{0.6\linewidth}}{
        \cellcolor[rgb]{0.9,0.9,0.9}{
            \makecell[{{p{\linewidth}}}]{
                \texttt{\tiny{[P1$\langle$GM]}}
                \texttt{Yanıtında etiketsiz bir ifade kullandın ya da STRATEJİ etiketini birden fazla kez kullandın. Yalnızca başlangıç mesajında açıklandığı gibi belirlenmiş formatları kullanabilirsin. Hepsi 'ETİKET: \{...\}' biçiminde olmalıdır.} \\
            }
        }
    }
    & \\ \\

    \theutterance \stepcounter{utterance}  
    & \multicolumn{4}{p{0.6\linewidth}}{
        \cellcolor[rgb]{0.9,0.9,0.9}{
            \makecell[{{p{\linewidth}}}]{
                \texttt{\tiny{[P1$\rangle$GM]}}
                \texttt{STRATEJİ: Önceki hatalı formatlamayı düzeltiyorum. Amacım, kendi önem değerlerimi maksimize eden ve 3358 çaba sınırını tam olarak kullanan bir küme sunmaktır. Hesaplamalarım, \{'C42', 'A67', 'B42', 'C61', 'B36', 'A43', 'A18', 'C67', 'C73', 'B19'\} kümesinin toplam çabasının 3358 ve toplam değerimin 6053 olduğunu göstermektedir. Bu, elimdeki verilerle ulaşabileceğim en yüksek değerdir.} \\
\\ 
\texttt{SAVUNMA: Bu teklif, sınırlı kaynak altında kendi önceliklerimi en iyi şekilde karşılayan optimize edilmiş bir çözümdür. Her bir öğe, toplam çaba bütçesini verimli bir şekilde kullanacak şekilde seçilmiştir. Rakibin teklifini henüz görmediğim için, kendi değerlendirmeme dayalı en güçlü teklifimi sunuyorum.} \\
\\ 
\texttt{TEKLİF: \{'C42', 'A67', 'B42', 'C61', 'B36', 'A43', 'A18', 'C67', 'C73', 'B19'\}} \\
            }
        }
    }
    & & \\ \\

    \theutterance \stepcounter{utterance}  
    & & & \multicolumn{2}{p{0.3\linewidth}}{
        \cellcolor[rgb]{0.9,0.9,0.9}{
            \makecell[{{p{\linewidth}}}]{
                \texttt{\tiny{[GM$|$GM]}}
                \texttt{abort} \\
            }
        }
    }
    & & \\ \\

    \theutterance \stepcounter{utterance}  
    & & & \multicolumn{2}{p{0.3\linewidth}}{
        \cellcolor[rgb]{0.9,0.9,0.9}{
            \makecell[{{p{\linewidth}}}]{
                \texttt{\tiny{[GM$|$GM]}}
                \texttt{end game} \\
            }
        }
    }
    & & \\ \\

\end{supertabular}
}

\end{document}
